\documentclass[12pt]{article}
\usepackage[utf8]{inputenc}
\usepackage[T1]{fontenc}
\usepackage{hyperref}
\title{Maléku Jaíka Language Resources for K-12 Curriculum Development}
\author{Compiled by Claude Sonnet 5.5 (claude-sonnet-5-5) for Andr\'as Kornai}
\date{}
\begin{document}
\maketitle

\section{Introduction}
Maléku Jaíka (also Guatuso; ISO 639-3 gut) is a Chibchan language of northern Costa Rica, spoken in the Guatuso area (villages of Margarita, Tonjibe and El Sol). Speaker estimates vary (one source: 230 to 750). It is classed as definitely endangered. Secondary descriptions in search results state that the Ministry of Education adopted a Latin-script orthography proposed by Adolfo Constenla Uma\~na and that the language is used alongside Spanish in some primary schools and taught as a subject in a secondary school; I did not verify this against primary sources.

\section{Primary Educational Resources}
No primer or textbook was identified in search results. A 2011 article in Revista Electr\'onica Educare on teaching indigenous languages in Costa Rican schools and integrating Maléku into the curriculum: \url{https://doaj.org/article/1fb5b2f1cac545bfb466b15ba40319ae}.

\section{Linguistic and Reference Materials}
Wikipedia overview: \url{https://en.wikipedia.org/wiki/Mal%C3%A9ku_language}. Omniglot page on the language and writing: \url{https://www.omniglot.com/writing/maleku.htm}. Endangered Languages Project entry: \url{https://www.endangeredlanguages.com/node/39280}. No dictionary was found in the searches.

\section{Religious and Community Resources}
A library record, ``Maleku: A Culture in Danger of Extinction'': \url{https://kimuk.conare.ac.cr/Record/PUNA_9e2859dc8fdceaa8df2f0fb7947e5693}.

\section{Digital and Audio Resources}
None found.

\section{Curriculum Development Context}
Any curriculum should build on the existing official orthography and on the bilingual intercultural education framework of Costa Rica, and should start by locating Ministry of Education materials and Constenla's publications, which these searches did not surface.

\end{document}